% =============================================================================
% Chapter 12: Συλλογή Λυμένων Ασκήσεων
% Data Mining Study Guide - NTUA
% =============================================================================

\chapter{Συλλογή Λυμένων Ασκήσεων}
\label{ch:solved-exercises}

Αυτό το κεφάλαιο περιέχει μια ολοκληρωμένη συλλογή λυμένων ασκήσεων που καλύπτουν
όλη την ύλη του μαθήματος Data Mining. Κάθε άσκηση συνοδεύεται από αναλυτική
επεξήγηση βήμα-βήμα.

% =============================================================================
% SECTION 1: E-R Model
% =============================================================================
\section{E-R Μοντέλο και Σχεδιασμός Βάσεων Δεδομένων}

% -----------------------------------------------------------------------------
% Exercise 1.1
% -----------------------------------------------------------------------------
\subsection{Άσκηση 1.1: Σχεδιασμός E-R από Απαιτήσεις}

\begin{formulabox}[Εκφώνηση]
Σχεδιάστε ένα E-R διάγραμμα για ένα σύστημα music streaming με τις εξής απαιτήσεις:
\begin{itemize}
    \item Κάθε \textbf{τραγούδι} έχει μοναδικό ID, τίτλο, διάρκεια και έτος κυκλοφορίας
    \item Κάθε \textbf{καλλιτέχνης} έχει μοναδικό ID, όνομα και χώρα προέλευσης
    \item Κάθε τραγούδι ανήκει σε ακριβώς έναν καλλιτέχνη, αλλά ένας καλλιτέχνης μπορεί να έχει πολλά τραγούδια
    \item Οι \textbf{χρήστες} έχουν ID, username και email
    \item Οι χρήστες δημιουργούν \textbf{playlists} με όνομα και ημερομηνία δημιουργίας
    \item Μια playlist μπορεί να περιέχει πολλά τραγούδια και ένα τραγούδι μπορεί να είναι σε πολλές playlists
\end{itemize}
\end{formulabox}

\subsubsection*{Λύση}

\textbf{Βήμα 1: Αναγνώριση Οντοτήτων}

\textit{Γιατί:} Πρώτα εντοπίζουμε τα «αντικείμενα» του συστήματος που έχουν ανεξάρτητη ύπαρξη.

Οντότητες:
\begin{itemize}
    \item \texttt{Song} (τραγούδι)
    \item \texttt{Artist} (καλλιτέχνης)
    \item \texttt{User} (χρήστης)
    \item \texttt{Playlist}
\end{itemize}

\textbf{Βήμα 2: Αναγνώριση Γνωρισμάτων}

\begin{center}
\begin{tabular}{ll}
\toprule
\textbf{Οντότητα} & \textbf{Γνωρίσματα} \\
\midrule
Song & \underline{song\_id}, title, duration, release\_year \\
Artist & \underline{artist\_id}, name, country \\
User & \underline{user\_id}, username, email \\
Playlist & \underline{playlist\_id}, name, created\_date \\
\bottomrule
\end{tabular}
\end{center}

\textbf{Βήμα 3: Αναγνώριση Συσχετίσεων και Cardinality}

\begin{itemize}
    \item \texttt{Artist} --- \texttt{creates} --- \texttt{Song}: \textbf{1:N}
    \item \texttt{User} --- \texttt{owns} --- \texttt{Playlist}: \textbf{1:N}
    \item \texttt{Playlist} --- \texttt{contains} --- \texttt{Song}: \textbf{M:N}
\end{itemize}

\textbf{Βήμα 4: E-R Διάγραμμα}

\begin{center}
\begin{tikzpicture}[node distance=3cm,
    entity/.style={rectangle, draw=primary, fill=box-blue, minimum width=2cm, minimum height=1cm, font=\bfseries},
    relationship/.style={diamond, draw=secondary, fill=box-yellow, aspect=2, font=\small},
    attribute/.style={ellipse, draw=gray, fill=light-gray, font=\small\ttfamily, inner sep=2pt}]

    % Entities
    \node[entity] (artist) {Artist};
    \node[entity, right=4cm of artist] (song) {Song};
    \node[entity, below=3cm of artist] (user) {User};
    \node[entity, right=4cm of user] (playlist) {Playlist};

    % Relationships
    \node[relationship] (creates) at ($(artist)!0.5!(song)$) {creates};
    \node[relationship] (owns) at ($(user)!0.5!(playlist)$) {owns};
    \node[relationship] (contains) at ($(playlist)!0.5!(song)$) {contains};

    % Connections with cardinality
    \draw[thick] (artist) -- node[above] {1} (creates);
    \draw[thick] (creates) -- node[above] {N} (song);
    \draw[thick] (user) -- node[above] {1} (owns);
    \draw[thick] (owns) -- node[above] {N} (playlist);
    \draw[thick] (playlist) -- node[left] {M} (contains);
    \draw[thick] (contains) -- node[right] {N} (song);

\end{tikzpicture}
\end{center}

\begin{warningbox}[Συνήθες Λάθος]
Πολλοί φοιτητές ξεχνούν ότι η σχέση \texttt{contains} είναι M:N και χρειάζεται
ενδιάμεσο πίνακα στο σχεσιακό μοντέλο!
\end{warningbox}

\textbf{Τελική Απάντηση:} Το E-R διάγραμμα περιλαμβάνει 4 οντότητες και 3 συσχετίσεις
(1:N, 1:N, M:N).

% -----------------------------------------------------------------------------
% Exercise 1.2
% -----------------------------------------------------------------------------
\subsection{Άσκηση 1.2: Μετατροπή E-R σε Σχεσιακό (M:N)}

\begin{formulabox}[Εκφώνηση]
Μετατρέψτε την M:N σχέση \texttt{contains} μεταξύ \texttt{Playlist} και \texttt{Song}
σε σχεσιακούς πίνακες. Η σχέση έχει επιπλέον γνώρισμα \texttt{position} που δηλώνει
τη θέση του τραγουδιού στη playlist.
\end{formulabox}

\subsubsection*{Λύση}

\textbf{Βήμα 1: Κανόνας Μετατροπής M:N}

\textit{Γιατί:} Οι M:N σχέσεις δεν μπορούν να αναπαρασταθούν με ξένα κλειδιά σε έναν
από τους δύο πίνακες. Χρειάζεται ενδιάμεσος πίνακας.

\textbf{Βήμα 2: Δημιουργία Πινάκων}

\begin{lstlisting}[language=SQL]
-- Πίνακας Playlist
CREATE TABLE Playlist (
    playlist_id INT PRIMARY KEY,
    name VARCHAR(100),
    created_date DATE,
    user_id INT REFERENCES User(user_id)
);

-- Πίνακας Song
CREATE TABLE Song (
    song_id INT PRIMARY KEY,
    title VARCHAR(200),
    duration INT,
    release_year INT,
    artist_id INT REFERENCES Artist(artist_id)
);

-- Ενδιάμεσος πίνακας για M:N
CREATE TABLE Playlist_Song (
    playlist_id INT REFERENCES Playlist(playlist_id),
    song_id INT REFERENCES Song(song_id),
    position INT,  -- Γνώρισμα της σχέσης
    PRIMARY KEY (playlist_id, song_id)
);
\end{lstlisting}

\begin{infobox}[Σημαντικό]
Το πρωτεύον κλειδί του ενδιάμεσου πίνακα είναι το συνδυασμένο κλειδί
\texttt{(playlist\_id, song\_id)} που αποτρέπει διπλότυπες εγγραφές.
\end{infobox}

\textbf{Τελική Απάντηση:} Δημιουργούνται 3 πίνακες: \texttt{Playlist}, \texttt{Song},
και \texttt{Playlist\_Song} (ενδιάμεσος).

% -----------------------------------------------------------------------------
% Exercise 1.3
% -----------------------------------------------------------------------------
\subsection{Άσκηση 1.3: Weak Entity Μετατροπή}

\begin{formulabox}[Εκφώνηση]
Ένα σύστημα διαχείρισης παραγγελιών έχει:
\begin{itemize}
    \item \texttt{Order} με \underline{order\_id} και order\_date
    \item \texttt{OrderItem} (weak entity) με \underline{item\_number}, quantity, price
\end{itemize}
Το \texttt{OrderItem} εξαρτάται από το \texttt{Order} για την αναγνώρισή του.
Μετατρέψτε σε σχεσιακό μοντέλο.
\end{formulabox}

\subsubsection*{Λύση}

\textbf{Βήμα 1: Αναγνώριση Weak Entity}

\textit{Γιατί:} Μια weak entity δεν έχει αρκετά γνωρίσματα για να ταυτοποιηθεί μόνη
της. Το \texttt{item\_number} (π.χ. 1, 2, 3) επαναλαμβάνεται μεταξύ διαφορετικών
παραγγελιών.

\textbf{Βήμα 2: Σύνθετο Πρωτεύον Κλειδί}

Το πρωτεύον κλειδί της weak entity σχηματίζεται από:
\begin{itemize}
    \item Το πρωτεύον κλειδί της «ιδιοκτήτριας» οντότητας (owner entity)
    \item Το μερικό κλειδί (partial key) της weak entity
\end{itemize}

\textbf{Βήμα 3: Σχεσιακοί Πίνακες}

\begin{lstlisting}[language=SQL]
CREATE TABLE Order (
    order_id INT PRIMARY KEY,
    order_date DATE
);

CREATE TABLE OrderItem (
    order_id INT REFERENCES Order(order_id),
    item_number INT,
    quantity INT,
    price DECIMAL(10,2),
    PRIMARY KEY (order_id, item_number)  -- Σύνθετο κλειδί
);
\end{lstlisting}

\begin{warningbox}[Συνήθες Λάθος]
Μην ξεχνάτε: το ξένο κλειδί (\texttt{order\_id}) γίνεται \textbf{μέρος} του
πρωτεύοντος κλειδιού στη weak entity, όχι απλά ξένο κλειδί!
\end{warningbox}

\textbf{Τελική Απάντηση:} PK του OrderItem = \texttt{(order\_id, item\_number)}

% -----------------------------------------------------------------------------
% Exercise 1.4
% -----------------------------------------------------------------------------
\subsection{Άσκηση 1.4: Cardinality Constraints}

\begin{formulabox}[Εκφώνηση]
Για τις παρακάτω περιγραφές, προσδιορίστε το cardinality (1:1, 1:N, M:N):

\begin{enumerate}
    \item Ένας φοιτητής εγγράφεται σε πολλά μαθήματα, ένα μάθημα έχει πολλούς φοιτητές
    \item Κάθε υπάλληλος έχει ακριβώς έναν διευθυντή, ένας διευθυντής διευθύνει πολλούς υπαλλήλους
    \item Κάθε χώρα έχει μία πρωτεύουσα, κάθε πρωτεύουσα ανήκει σε μία χώρα
    \item Ένας συγγραφέας γράφει πολλά βιβλία, ένα βιβλίο μπορεί να έχει πολλούς συγγραφείς
\end{enumerate}
\end{formulabox}

\subsubsection*{Λύση}

\textbf{Μέθοδος:} Για κάθε περίπτωση ρωτάμε:
\begin{itemize}
    \item Πόσες οντότητες Β συνδέονται με μία οντότητα Α;
    \item Πόσες οντότητες Α συνδέονται με μία οντότητα Β;
\end{itemize}

\begin{center}
\begin{tabular}{clcl}
\toprule
\textbf{\#} & \textbf{Σχέση} & \textbf{Cardinality} & \textbf{Εξήγηση} \\
\midrule
1 & Student-Course & \textbf{M:N} & Πολλοί-σε-πολλούς \\
2 & Employee-Manager & \textbf{N:1} & Πολλοί υπάλληλοι, ένας διευθυντής \\
3 & Country-Capital & \textbf{1:1} & Ένα-προς-ένα \\
4 & Author-Book & \textbf{M:N} & Πολλοί-σε-πολλούς \\
\bottomrule
\end{tabular}
\end{center}

\begin{infobox}[Υλοποίηση σε Σχεσιακό]
\begin{itemize}
    \item \textbf{1:1}: FK σε οποιονδήποτε πίνακα (+ UNIQUE)
    \item \textbf{1:N}: FK στην πλευρά του N
    \item \textbf{M:N}: Ενδιάμεσος πίνακας (junction table)
\end{itemize}
\end{infobox}

% -----------------------------------------------------------------------------
% Exercise 1.5
% -----------------------------------------------------------------------------
\subsection{Άσκηση 1.5: Σύνθετο Schema (Hospital System)}

\begin{formulabox}[Εκφώνηση]
Σχεδιάστε E-R και μετατρέψτε σε σχεσιακό για νοσοκομειακό σύστημα:
\begin{itemize}
    \item \texttt{Doctor}: doctor\_id, name, specialty
    \item \texttt{Patient}: patient\_id, name, birth\_date
    \item \texttt{Appointment}: date, time, diagnosis (σχέση μεταξύ Doctor-Patient)
    \item Ένας γιατρός έχει πολλά ραντεβού, ένας ασθενής έχει πολλά ραντεβού
    \item Ένα ραντεβού συνδέει ακριβώς έναν γιατρό με έναν ασθενή
\end{itemize}
\end{formulabox}

\subsubsection*{Λύση}

\textbf{Βήμα 1: Ανάλυση}

Η σχέση Doctor-Patient είναι M:N (ένας γιατρός βλέπει πολλούς ασθενείς, ένας ασθενής
επισκέπτεται πολλούς γιατρούς), αλλά κάθε \textit{συγκεκριμένο} ραντεβού αφορά
ακριβώς έναν γιατρό και έναν ασθενή.

\textbf{Βήμα 2: E-R Διάγραμμα}

\begin{center}
\begin{tikzpicture}[node distance=4cm,
    entity/.style={rectangle, draw=primary, fill=box-blue, minimum width=2cm, minimum height=1cm, font=\bfseries},
    relationship/.style={diamond, draw=secondary, fill=box-yellow, aspect=2.5, font=\small}]

    \node[entity] (doctor) {Doctor};
    \node[entity, right=5cm of doctor] (patient) {Patient};
    \node[relationship] (appointment) at ($(doctor)!0.5!(patient)$) {Appointment};

    \draw[thick] (doctor) -- node[above] {1} (appointment);
    \draw[thick] (appointment) -- node[above] {N} (patient);
    \draw[thick] (patient) -- node[below] {1} (appointment);
    \draw[thick] (appointment) -- node[below] {M} (doctor);

    % Attributes of relationship
    \node[ellipse, draw=gray, fill=light-gray, font=\small, above=0.5cm of appointment] (date) {date, time};
    \node[ellipse, draw=gray, fill=light-gray, font=\small, below=0.5cm of appointment] (diag) {diagnosis};
    \draw[dashed] (appointment) -- (date);
    \draw[dashed] (appointment) -- (diag);
\end{tikzpicture}
\end{center}

\textbf{Βήμα 3: Σχεσιακό Μοντέλο}

\begin{lstlisting}[language=SQL]
CREATE TABLE Doctor (
    doctor_id INT PRIMARY KEY,
    name VARCHAR(100),
    specialty VARCHAR(50)
);

CREATE TABLE Patient (
    patient_id INT PRIMARY KEY,
    name VARCHAR(100),
    birth_date DATE
);

CREATE TABLE Appointment (
    appointment_id INT PRIMARY KEY,
    doctor_id INT REFERENCES Doctor(doctor_id),
    patient_id INT REFERENCES Patient(patient_id),
    appointment_date DATE,
    appointment_time TIME,
    diagnosis TEXT
);
\end{lstlisting}

\textbf{Τελική Απάντηση:} 3 πίνακες με τον \texttt{Appointment} να περιέχει FKs και
προς τους δύο.


% =============================================================================
% SECTION 2: Relational Algebra
% =============================================================================
\section{Σχεσιακή Άλγεβρα}

\begin{infobox}[Σύνοψη Τελεστών]
\begin{center}
\begin{tabular}{cll}
\toprule
\textbf{Σύμβολο} & \textbf{Όνομα} & \textbf{Περιγραφή} \\
\midrule
$\sigma$ & Selection & Επιλογή γραμμών \\
$\pi$ & Projection & Επιλογή στηλών \\
$\cup$ & Union & Ένωση συνόλων \\
$\cap$ & Intersection & Τομή συνόλων \\
$-$ & Difference & Διαφορά συνόλων \\
$\times$ & Cartesian Product & Καρτεσιανό γινόμενο \\
$\bowtie$ & Natural Join & Φυσική σύζευξη \\
$\bowtie_\theta$ & Theta Join & Σύζευξη με συνθήκη \\
$\div$ & Division & Διαίρεση \\
\bottomrule
\end{tabular}
\end{center}
\end{infobox}

% -----------------------------------------------------------------------------
% Exercise 2.1
% -----------------------------------------------------------------------------
\subsection{Άσκηση 2.1: Selection και Projection}

\begin{formulabox}[Εκφώνηση]
Δίνεται ο πίνακας \texttt{Employee(emp\_id, name, department, salary)}.

Γράψτε σε σχεσιακή άλγεβρα: «Βρες τα ονόματα των υπαλλήλων του τμήματος IT με
μισθό πάνω από 50000».
\end{formulabox}

\subsubsection*{Λύση}

\textbf{Βήμα 1: Selection (επιλογή γραμμών)}

\textit{Γιατί:} Πρώτα φιλτράρουμε τις γραμμές που πληρούν τις συνθήκες.

$$\sigma_{department='IT' \land salary > 50000}(Employee)$$

\textbf{Βήμα 2: Projection (επιλογή στηλών)}

\textit{Γιατί:} Μετά κρατάμε μόνο τη στήλη που μας ενδιαφέρει.

$$\pi_{name}(\sigma_{department='IT' \land salary > 50000}(Employee))$$

\begin{warningbox}[Σειρά Εφαρμογής]
Η σειρά Selection $\to$ Projection είναι σημαντική για βελτιστοποίηση! Αν κάνουμε
πρώτα projection, χάνουμε τις στήλες που χρειαζόμαστε για το φιλτράρισμα.
\end{warningbox}

\textbf{Τελική Απάντηση:} $\pi_{name}(\sigma_{department='IT' \land salary > 50000}(Employee))$

% -----------------------------------------------------------------------------
% Exercise 2.2
% -----------------------------------------------------------------------------
\subsection{Άσκηση 2.2: Natural Join}

\begin{formulabox}[Εκφώνηση]
Δίνονται:
\begin{itemize}
    \item \texttt{Student(sid, sname, dept)}
    \item \texttt{Enrollment(sid, cid, grade)}
    \item \texttt{Course(cid, cname, credits)}
\end{itemize}

Βρείτε τα ονόματα φοιτητών και τα μαθήματα στα οποία είναι εγγεγραμμένοι.
\end{formulabox}

\subsubsection*{Λύση}

\textbf{Βήμα 1: Αναγνώριση κοινών γνωρισμάτων}

\begin{itemize}
    \item Student $\bowtie$ Enrollment: κοινό \texttt{sid}
    \item Enrollment $\bowtie$ Course: κοινό \texttt{cid}
\end{itemize}

\textbf{Βήμα 2: Natural Join αλυσίδα}

$$Student \bowtie Enrollment \bowtie Course$$

\textbf{Βήμα 3: Projection}

$$\pi_{sname, cname}(Student \bowtie Enrollment \bowtie Course)$$

\begin{infobox}[Πώς λειτουργεί το Natural Join]
Το $\bowtie$ συνδέει πλειάδες που έχουν \textbf{ίδιες τιμές} στα κοινά γνωρίσματα
(αυτά με το ίδιο όνομα). Οι κοινές στήλες εμφανίζονται μία φορά στο αποτέλεσμα.
\end{infobox}

% -----------------------------------------------------------------------------
% Exercise 2.3
% -----------------------------------------------------------------------------
\subsection{Άσκηση 2.3: Theta Join}

\begin{formulabox}[Εκφώνηση]
Δίνονται:
\begin{itemize}
    \item \texttt{Employee(eid, ename, salary, manager\_id)}
    \item \texttt{Manager(mid, mname, bonus)}
\end{itemize}

Βρείτε ζεύγη (υπάλληλος, διευθυντής) όπου ο μισθός του υπαλλήλου είναι
μεγαλύτερος από το bonus του διευθυντή του.
\end{formulabox}

\subsubsection*{Λύση}

\textbf{Βήμα 1: Σύνδεση υπαλλήλου με διευθυντή}

Πρώτα συνδέουμε με βάση το \texttt{manager\_id = mid}:
$$Employee \bowtie_{manager\_id = mid} Manager$$

\textbf{Βήμα 2: Φιλτράρισμα}

Εφαρμόζουμε τη συνθήκη \texttt{salary > bonus}:
$$\sigma_{salary > bonus}(Employee \bowtie_{manager\_id = mid} Manager)$$

\textbf{Βήμα 3: Projection}

$$\pi_{ename, mname}(\sigma_{salary > bonus}(Employee \bowtie_{manager\_id = mid} Manager))$$

\textbf{Εναλλακτικά (ενιαίο Theta Join):}
$$\pi_{ename, mname}(Employee \bowtie_{manager\_id = mid \land salary > bonus} Manager)$$

% -----------------------------------------------------------------------------
% Exercise 2.4
% -----------------------------------------------------------------------------
\subsection{Άσκηση 2.4: Set Difference}

\begin{formulabox}[Εκφώνηση]
Δίνονται:
\begin{itemize}
    \item \texttt{AllStudents(sid, sname)}
    \item \texttt{Enrollment(sid, cid)}
\end{itemize}

Βρείτε τους φοιτητές που \textbf{δεν} είναι εγγεγραμμένοι σε κανένα μάθημα.
\end{formulabox}

\subsubsection*{Λύση}

\textbf{Βήμα 1: Εύρεση εγγεγραμμένων φοιτητών}

$$EnrolledStudents = \pi_{sid}(Enrollment)$$

\textbf{Βήμα 2: Set Difference}

\textit{Γιατί:} Το $A - B$ δίνει τα στοιχεία που είναι στο A αλλά όχι στο B.

$$NotEnrolled = \pi_{sid}(AllStudents) - \pi_{sid}(Enrollment)$$

\textbf{Βήμα 3: Επαναφορά ονομάτων (αν χρειάζεται)}

$$\pi_{sname}(AllStudents \bowtie NotEnrolled)$$

\begin{warningbox}[Απαίτηση Συμβατότητας]
Για τελεστές συνόλων ($\cup$, $\cap$, $-$), οι σχέσεις πρέπει να έχουν:
\begin{enumerate}
    \item Ίδιο αριθμό γνωρισμάτων
    \item Συμβατά domains στις αντίστοιχες θέσεις
\end{enumerate}
\end{warningbox}

% -----------------------------------------------------------------------------
% Exercise 2.5
% -----------------------------------------------------------------------------
\subsection{Άσκηση 2.5: Intersection με Projection}

\begin{formulabox}[Εκφώνηση]
Δίνονται:
\begin{itemize}
    \item \texttt{CS\_Students(sid, sname)} -- Φοιτητές πληροφορικής
    \item \texttt{Math\_Students(sid, sname)} -- Φοιτητές μαθηματικών
\end{itemize}

Βρείτε τους φοιτητές που σπουδάζουν \textbf{και} πληροφορική \textbf{και} μαθηματικά.
\end{formulabox}

\subsubsection*{Λύση}

$$CS\_Students \cap Math\_Students$$

ή ισοδύναμα (αν θέλουμε μόνο ονόματα):

$$\pi_{sname}(CS\_Students \cap Math\_Students)$$

\begin{infobox}[Ισοδυναμία]
Η τομή μπορεί να εκφραστεί με διαφορά:
$$A \cap B = A - (A - B)$$
\end{infobox}

% -----------------------------------------------------------------------------
% Exercise 2.6
% -----------------------------------------------------------------------------
\subsection{Άσκηση 2.6: Division (Απλή)}

\begin{formulabox}[Εκφώνηση]
Δίνονται:
\begin{itemize}
    \item \texttt{Takes(sid, cid)} -- φοιτητής παίρνει μάθημα
    \item \texttt{Required(cid)} -- υποχρεωτικά μαθήματα: $\{C1, C2\}$
\end{itemize}

Βρείτε τους φοιτητές που έχουν πάρει \textbf{όλα} τα υποχρεωτικά μαθήματα.
\end{formulabox}

\subsubsection*{Λύση}

\textbf{Η Division ($\div$) απαντά σε ερωτήματα «για όλα»:}

$$Takes \div Required$$

\textbf{Βήμα-βήμα αλγόριθμος:}

\begin{enumerate}
    \item Βρες όλους τους φοιτητές: $\pi_{sid}(Takes)$
    \item Βρες τι «θα έπρεπε» να έχει πάρει κάθε φοιτητής: $\pi_{sid}(Takes) \times Required$
    \item Βρες τι «δεν έχει» πάρει: $(\pi_{sid}(Takes) \times Required) - Takes$
    \item Αφαίρεσε όσους «δεν έχουν» κάτι: $\pi_{sid}(Takes) - \pi_{sid}((\pi_{sid}(Takes) \times Required) - Takes)$
\end{enumerate}

\textbf{Παράδειγμα:}

\begin{center}
\begin{tabular}{cc}
\textbf{Takes} & \\
\toprule
sid & cid \\
\midrule
S1 & C1 \\
S1 & C2 \\
S1 & C3 \\
S2 & C1 \\
S3 & C1 \\
S3 & C2 \\
\bottomrule
\end{tabular}
\quad\quad
\begin{tabular}{c}
\textbf{Required} \\
\toprule
cid \\
\midrule
C1 \\
C2 \\
\bottomrule
\end{tabular}
\quad\quad
\begin{tabular}{c}
\textbf{Result} \\
\toprule
sid \\
\midrule
S1 \\
S3 \\
\bottomrule
\end{tabular}
\end{center}

% -----------------------------------------------------------------------------
% Exercise 2.7
% -----------------------------------------------------------------------------
\subsection{Άσκηση 2.7: Division (Σύνθετη)}

\begin{formulabox}[Εκφώνηση]
Δίνονται:
\begin{itemize}
    \item \texttt{Supplies(supplier, part)}
    \item \texttt{Project(pname, part)} -- έργα και τα parts που χρειάζονται
\end{itemize}

Βρείτε τους προμηθευτές που μπορούν να προμηθεύσουν \textbf{όλα} τα parts που
χρειάζεται το έργο «Rocket».
\end{formulabox}

\subsubsection*{Λύση}

\textbf{Βήμα 1: Βρες τα parts του Rocket}

$$RocketParts = \pi_{part}(\sigma_{pname='Rocket'}(Project))$$

\textbf{Βήμα 2: Division}

$$Supplies \div RocketParts$$

\textbf{Σε SQL (για κατανόηση):}

\begin{lstlisting}[language=SQL]
SELECT DISTINCT s1.supplier
FROM Supplies s1
WHERE NOT EXISTS (
    SELECT part FROM Project WHERE pname = 'Rocket'
    EXCEPT
    SELECT part FROM Supplies s2
    WHERE s2.supplier = s1.supplier
);
\end{lstlisting}

% -----------------------------------------------------------------------------
% Exercise 2.8
% -----------------------------------------------------------------------------
\subsection{Άσκηση 2.8: Σύνθετη Έκφραση}

\begin{formulabox}[Εκφώνηση]
Δίνονται:
\begin{itemize}
    \item \texttt{Employee(eid, ename, dept\_id, salary)}
    \item \texttt{Department(dept\_id, dname, budget)}
    \item \texttt{Project(pid, pname, dept\_id)}
    \item \texttt{WorksOn(eid, pid, hours)}
\end{itemize}

Βρείτε τα ονόματα υπαλλήλων του τμήματος «Research» που εργάζονται πάνω από 20 ώρες
σε κάποιο project.
\end{formulabox}

\subsubsection*{Λύση}

\textbf{Βήμα 1: Εύρεση dept\_id του Research}

$$ResearchDept = \sigma_{dname='Research'}(Department)$$

\textbf{Βήμα 2: Υπάλληλοι του Research}

$$ResearchEmps = Employee \bowtie ResearchDept$$

\textbf{Βήμα 3: Υπάλληλοι με >20 ώρες}

$$HighHours = \sigma_{hours > 20}(WorksOn)$$

\textbf{Βήμα 4: Συνδυασμός}

$$Result = \pi_{ename}(ResearchEmps \bowtie HighHours)$$

\textbf{Ενιαία έκφραση:}

$$\pi_{ename}(\sigma_{dname='Research'}(Department) \bowtie Employee \bowtie \sigma_{hours>20}(WorksOn))$$


% =============================================================================
% SECTION 3: SQL
% =============================================================================
\section{SQL}

% -----------------------------------------------------------------------------
% Exercise 3.1
% -----------------------------------------------------------------------------
\subsection{Άσκηση 3.1: SELECT με WHERE}

\begin{formulabox}[Εκφώνηση]
Δίνεται ο πίνακας \texttt{Products(pid, pname, category, price, stock)}.

Γράψτε SQL query για να βρείτε τα προϊόντα κατηγορίας «Electronics» με τιμή μεταξύ
100 και 500 ευρώ, ταξινομημένα κατά τιμή (φθίνουσα).
\end{formulabox}

\subsubsection*{Λύση}

\begin{lstlisting}[language=SQL]
SELECT pid, pname, price
FROM Products
WHERE category = 'Electronics'
  AND price BETWEEN 100 AND 500
ORDER BY price DESC;
\end{lstlisting}

\begin{infobox}[BETWEEN]
Το \texttt{BETWEEN 100 AND 500} είναι ισοδύναμο με \texttt{price >= 100 AND price <= 500}
(περιλαμβάνει τα όρια).
\end{infobox}

% -----------------------------------------------------------------------------
% Exercise 3.2
% -----------------------------------------------------------------------------
\subsection{Άσκηση 3.2: JOINs (INNER, LEFT)}

\begin{formulabox}[Εκφώνηση]
Δίνονται:
\begin{itemize}
    \item \texttt{Customers(cid, cname, city)}
    \item \texttt{Orders(oid, cid, order\_date, total)}
\end{itemize}

\begin{enumerate}
    \item Βρείτε όλους τους πελάτες με τις παραγγελίες τους
    \item Βρείτε όλους τους πελάτες, ακόμα κι αν δεν έχουν παραγγελίες
\end{enumerate}
\end{formulabox}

\subsubsection*{Λύση}

\textbf{(a) INNER JOIN -- μόνο πελάτες με παραγγελίες:}

\begin{lstlisting}[language=SQL]
SELECT c.cname, o.oid, o.total
FROM Customers c
INNER JOIN Orders o ON c.cid = o.cid;
\end{lstlisting}

\textbf{(b) LEFT JOIN -- όλοι οι πελάτες:}

\begin{lstlisting}[language=SQL]
SELECT c.cname, o.oid, o.total
FROM Customers c
LEFT JOIN Orders o ON c.cid = o.cid;
\end{lstlisting}

\begin{center}
\begin{tikzpicture}[scale=0.8]
    % INNER JOIN
    \begin{scope}[xshift=-4cm]
        \draw[fill=box-blue, fill opacity=0.5] (0,0) circle (1.2);
        \draw[fill=box-green, fill opacity=0.5] (1.5,0) circle (1.2);
        \fill[secondary, opacity=0.7] (0.75,0) circle (0.5);
        \node at (0.75,-2) {\small INNER JOIN};
    \end{scope}
    % LEFT JOIN
    \begin{scope}[xshift=4cm]
        \draw[fill=secondary, fill opacity=0.5] (0,0) circle (1.2);
        \draw[fill=box-green, fill opacity=0.3] (1.5,0) circle (1.2);
        \node at (0.75,-2) {\small LEFT JOIN};
    \end{scope}
\end{tikzpicture}
\end{center}

% -----------------------------------------------------------------------------
% Exercise 3.3
% -----------------------------------------------------------------------------
\subsection{Άσκηση 3.3: GROUP BY με COUNT}

\begin{formulabox}[Εκφώνηση]
Χρησιμοποιώντας τους πίνακες \texttt{Orders} και \texttt{Customers}, βρείτε πόσες
παραγγελίες έχει κάνει κάθε πελάτης.
\end{formulabox}

\subsubsection*{Λύση}

\begin{lstlisting}[language=SQL]
SELECT c.cid, c.cname, COUNT(o.oid) AS order_count
FROM Customers c
LEFT JOIN Orders o ON c.cid = o.cid
GROUP BY c.cid, c.cname;
\end{lstlisting}

\begin{warningbox}[LEFT vs INNER JOIN]
Χρησιμοποιούμε LEFT JOIN για να συμπεριλάβουμε πελάτες με 0 παραγγελίες.
Με INNER JOIN θα χάναμε αυτούς τους πελάτες!
\end{warningbox}

% -----------------------------------------------------------------------------
% Exercise 3.4
% -----------------------------------------------------------------------------
\subsection{Άσκηση 3.4: GROUP BY με HAVING}

\begin{formulabox}[Εκφώνηση]
Βρείτε τις πόλεις που έχουν περισσότερους από 5 πελάτες και εμφανίστε τον αριθμό
πελατών ανά πόλη.
\end{formulabox}

\subsubsection*{Λύση}

\begin{lstlisting}[language=SQL]
SELECT city, COUNT(*) AS customer_count
FROM Customers
GROUP BY city
HAVING COUNT(*) > 5
ORDER BY customer_count DESC;
\end{lstlisting}

\begin{infobox}[WHERE vs HAVING]
\begin{itemize}
    \item \textbf{WHERE}: Φιλτράρει \textit{γραμμές} \textbf{πριν} το grouping
    \item \textbf{HAVING}: Φιλτράρει \textit{groups} \textbf{μετά} το grouping
\end{itemize}
\end{infobox}

% -----------------------------------------------------------------------------
% Exercise 3.5
% -----------------------------------------------------------------------------
\subsection{Άσκηση 3.5: Subquery με IN}

\begin{formulabox}[Εκφώνηση]
Βρείτε τους πελάτες που έχουν κάνει παραγγελία με total > 1000.
\end{formulabox}

\subsubsection*{Λύση}

\textbf{Μέθοδος 1: Subquery με IN}

\begin{lstlisting}[language=SQL]
SELECT cname
FROM Customers
WHERE cid IN (
    SELECT cid
    FROM Orders
    WHERE total > 1000
);
\end{lstlisting}

\textbf{Μέθοδος 2: JOIN (ισοδύναμο)}

\begin{lstlisting}[language=SQL]
SELECT DISTINCT c.cname
FROM Customers c
JOIN Orders o ON c.cid = o.cid
WHERE o.total > 1000;
\end{lstlisting}

% -----------------------------------------------------------------------------
% Exercise 3.6
% -----------------------------------------------------------------------------
\subsection{Άσκηση 3.6: Correlated Subquery}

\begin{formulabox}[Εκφώνηση]
Βρείτε τους υπαλλήλους που έχουν μισθό μεγαλύτερο από τον μέσο όρο του τμήματός τους.
\end{formulabox}

\subsubsection*{Λύση}

\begin{lstlisting}[language=SQL]
SELECT e1.ename, e1.salary, e1.dept_id
FROM Employee e1
WHERE e1.salary > (
    SELECT AVG(e2.salary)
    FROM Employee e2
    WHERE e2.dept_id = e1.dept_id  -- Correlated!
);
\end{lstlisting}

\begin{infobox}[Correlated Subquery]
Το subquery εξαρτάται από το εξωτερικό query (αναφέρεται στο \texttt{e1.dept\_id}).
Εκτελείται μία φορά \textbf{για κάθε γραμμή} του εξωτερικού query.
\end{infobox}

% -----------------------------------------------------------------------------
% Exercise 3.7
% -----------------------------------------------------------------------------
\subsection{Άσκηση 3.7: NOT EXISTS Pattern}

\begin{formulabox}[Εκφώνηση]
Βρείτε τους πελάτες που \textbf{δεν} έχουν κάνει καμία παραγγελία.
\end{formulabox}

\subsubsection*{Λύση}

\textbf{Μέθοδος 1: NOT EXISTS}

\begin{lstlisting}[language=SQL]
SELECT c.cname
FROM Customers c
WHERE NOT EXISTS (
    SELECT 1
    FROM Orders o
    WHERE o.cid = c.cid
);
\end{lstlisting}

\textbf{Μέθοδος 2: LEFT JOIN + IS NULL}

\begin{lstlisting}[language=SQL]
SELECT c.cname
FROM Customers c
LEFT JOIN Orders o ON c.cid = o.cid
WHERE o.oid IS NULL;
\end{lstlisting}

\textbf{Μέθοδος 3: NOT IN}

\begin{lstlisting}[language=SQL]
SELECT cname
FROM Customers
WHERE cid NOT IN (SELECT cid FROM Orders);
\end{lstlisting}

\begin{warningbox}[NOT IN με NULL]
Αν το subquery του NOT IN επιστρέψει NULL, το αποτέλεσμα είναι κενό!
Προτιμήστε NOT EXISTS για ασφάλεια.
\end{warningbox}

% -----------------------------------------------------------------------------
% Exercise 3.8
% -----------------------------------------------------------------------------
\subsection{Άσκηση 3.8: SQL Division (Double NOT EXISTS)}

\begin{formulabox}[Εκφώνηση]
Δίνονται:
\begin{itemize}
    \item \texttt{Enrollment(sid, cid)}
    \item \texttt{Course(cid, cname)}
\end{itemize}

Βρείτε τους φοιτητές που έχουν εγγραφεί σε \textbf{όλα} τα μαθήματα.
\end{formulabox}

\subsubsection*{Λύση}

\textbf{Λογική:} «Δεν υπάρχει μάθημα στο οποίο ο φοιτητής να μην έχει εγγραφεί»

\begin{lstlisting}[language=SQL]
SELECT DISTINCT e1.sid
FROM Enrollment e1
WHERE NOT EXISTS (
    -- Μαθήματα που ΔΕΝ έχει πάρει ο e1.sid
    SELECT c.cid
    FROM Course c
    WHERE NOT EXISTS (
        -- Έλεγχος αν ο e1.sid έχει πάρει το c.cid
        SELECT 1
        FROM Enrollment e2
        WHERE e2.sid = e1.sid AND e2.cid = c.cid
    )
);
\end{lstlisting}

\begin{infobox}[Ερμηνεία Double NOT EXISTS]
\begin{enumerate}
    \item Εσωτερικό NOT EXISTS: Βρες μαθήματα που \textbf{δεν} έχει πάρει ο φοιτητής
    \item Εξωτερικό NOT EXISTS: Κράτα φοιτητές που \textbf{δεν} έχουν τέτοια μαθήματα
\end{enumerate}
Δηλαδή: φοιτητές χωρίς «κενά» στο πρόγραμμα.
\end{infobox}

% -----------------------------------------------------------------------------
% Exercise 3.9
% -----------------------------------------------------------------------------
\subsection{Άσκηση 3.9: Multiple JOINs με Aggregation}

\begin{formulabox}[Εκφώνηση]
Δίνονται:
\begin{itemize}
    \item \texttt{Orders(oid, cid, order\_date)}
    \item \texttt{OrderItems(oid, pid, quantity, unit\_price)}
    \item \texttt{Products(pid, pname, category)}
\end{itemize}

Βρείτε το συνολικό έσοδο (revenue) ανά κατηγορία προϊόντος.
\end{formulabox}

\subsubsection*{Λύση}

\begin{lstlisting}[language=SQL]
SELECT
    p.category,
    SUM(oi.quantity * oi.unit_price) AS total_revenue
FROM Products p
JOIN OrderItems oi ON p.pid = oi.pid
GROUP BY p.category
ORDER BY total_revenue DESC;
\end{lstlisting}

% -----------------------------------------------------------------------------
% Exercise 3.10
% -----------------------------------------------------------------------------
\subsection{Άσκηση 3.10: Complex Query}

\begin{formulabox}[Εκφώνηση]
Βρείτε τους top 3 πελάτες (με βάση συνολικές αγορές) που έχουν αγοράσει προϊόντα
από τουλάχιστον 2 διαφορετικές κατηγορίες.
\end{formulabox}

\subsubsection*{Λύση}

\begin{lstlisting}[language=SQL]
SELECT
    c.cid,
    c.cname,
    SUM(oi.quantity * oi.unit_price) AS total_spent,
    COUNT(DISTINCT p.category) AS categories_bought
FROM Customers c
JOIN Orders o ON c.cid = o.cid
JOIN OrderItems oi ON o.oid = oi.oid
JOIN Products p ON oi.pid = p.pid
GROUP BY c.cid, c.cname
HAVING COUNT(DISTINCT p.category) >= 2
ORDER BY total_spent DESC
LIMIT 3;
\end{lstlisting}

\begin{infobox}[Σειρά Εκτέλεσης SQL]
\begin{enumerate}
    \item FROM / JOIN
    \item WHERE
    \item GROUP BY
    \item HAVING
    \item SELECT
    \item ORDER BY
    \item LIMIT
\end{enumerate}
\end{infobox}


% =============================================================================
% SECTION 4: Query Processing
% =============================================================================
\section{Query Processing}

% -----------------------------------------------------------------------------
% Exercise 4.1
% -----------------------------------------------------------------------------
\subsection{Άσκηση 4.1: Pipelining vs Materialization}

\begin{formulabox}[Εκφώνηση]
Για το query: $\pi_{name}(\sigma_{salary>50000}(Employee \bowtie Department))$

Ταξινομήστε τους τελεστές σε pipelined ή materialized και εξηγήστε γιατί.
\end{formulabox}

\subsubsection*{Λύση}

\textbf{Ορισμοί:}
\begin{itemize}
    \item \textbf{Pipelining:} Τα αποτελέσματα περνάνε αμέσως στον επόμενο τελεστή χωρίς αποθήκευση
    \item \textbf{Materialization:} Τα αποτελέσματα αποθηκεύονται προσωρινά πριν περάσουν στον επόμενο τελεστή
\end{itemize}

\begin{center}
\begin{tabular}{lll}
\toprule
\textbf{Τελεστής} & \textbf{Τύπος} & \textbf{Λόγος} \\
\midrule
$\bowtie$ (Join) & Materialized & Χρειάζεται πλήρη δεδομένα για hash/sort \\
$\sigma$ (Selection) & Pipelined & Φιλτράρει γραμμή-γραμμή \\
$\pi$ (Projection) & Pipelined & Επεξεργάζεται γραμμή-γραμμή \\
\bottomrule
\end{tabular}
\end{center}

\begin{infobox}[Pipeline Breakers]
Οι τελεστές που \textbf{δεν} μπορούν να είναι pipelined ονομάζονται «pipeline breakers»:
\begin{itemize}
    \item Sort
    \item Hash Join (build phase)
    \item Group By / Aggregation
    \item Duplicate elimination (DISTINCT)
\end{itemize}
\end{infobox}

% -----------------------------------------------------------------------------
% Exercise 4.2
% -----------------------------------------------------------------------------
\subsection{Άσκηση 4.2: Επιλογή Join Algorithm}

\begin{formulabox}[Εκφώνηση]
Δίνονται δύο πίνακες:
\begin{itemize}
    \item $R$: 10,000 tuples, 100 tuples/page $\Rightarrow$ 100 pages
    \item $S$: 1,000 tuples, 100 tuples/page $\Rightarrow$ 10 pages
    \item Buffer: 12 pages
\end{itemize}

Ποιος αλγόριθμος join είναι καλύτερος; Υπολογίστε το I/O cost.
\end{formulabox}

\subsubsection*{Λύση}

\textbf{Nested Loop Join:}

$$Cost_{NLJ} = B_R + B_R \cdot B_S = 100 + 100 \cdot 10 = 1100 \text{ I/Os}$$

(με το μικρότερο ως outer: $10 + 10 \cdot 100 = 1010$ I/Os)

\textbf{Block Nested Loop Join:}

Με buffer $B = 12$ pages, κρατάμε $B-2 = 10$ pages για τον outer.

$$Cost_{BNLJ} = B_R + \lceil B_R / (B-2) \rceil \cdot B_S$$

Με $S$ ως outer (μικρότερο):
$$Cost = 10 + \lceil 10/10 \rceil \cdot 100 = 10 + 100 = 110 \text{ I/Os}$$

\textbf{Hash Join:}

$$Cost_{Hash} = 3 \cdot (B_R + B_S) = 3 \cdot 110 = 330 \text{ I/Os}$$

\textbf{Sort-Merge Join:}

$$Cost_{SM} = Sort(R) + Sort(S) + B_R + B_S$$
$$= 2 \cdot 100 \cdot \lceil \log_{11}(10) \rceil + 2 \cdot 10 \cdot \lceil \log_{11}(1) \rceil + 110$$
$$\approx 200 \cdot 1 + 20 \cdot 1 + 110 = 330 \text{ I/Os}$$

\begin{center}
\begin{tabular}{lr}
\toprule
\textbf{Algorithm} & \textbf{I/O Cost} \\
\midrule
Simple NLJ (R outer) & 1,100 \\
Simple NLJ (S outer) & 1,010 \\
Block NLJ (S outer) & \textbf{110} \\
Hash Join & 330 \\
Sort-Merge & $\sim$330 \\
\bottomrule
\end{tabular}
\end{center}

\textbf{Τελική Απάντηση:} Block Nested Loop Join με S ως outer (110 I/Os).

% -----------------------------------------------------------------------------
% Exercise 4.3
% -----------------------------------------------------------------------------
\subsection{Άσκηση 4.3: Query Optimization (Push Selection)}

\begin{formulabox}[Εκφώνηση]
Βελτιστοποιήστε το query tree για:

$$\pi_{ename}(\sigma_{dname='Sales'}(Employee \bowtie Department))$$
\end{formulabox}

\subsubsection*{Λύση}

\textbf{Αρχικό Query Tree:}

\begin{center}
\begin{tikzpicture}[
    node distance=1.2cm,
    op/.style={circle, draw=secondary, fill=box-blue, minimum size=0.8cm},
    rel/.style={rectangle, draw=primary, fill=box-green, rounded corners, minimum width=1.5cm}]

    \node[op] (proj) {$\pi$};
    \node[op, below=of proj] (sel) {$\sigma$};
    \node[op, below=of sel] (join) {$\bowtie$};
    \node[rel, below left=1cm and 0.5cm of join] (emp) {Employee};
    \node[rel, below right=1cm and 0.5cm of join] (dept) {Department};

    \draw[arrow] (sel) -- (proj);
    \draw[arrow] (join) -- (sel);
    \draw[arrow] (emp) -- (join);
    \draw[arrow] (dept) -- (join);

    \node[right=0.3cm of proj, font=\small] {$name$};
    \node[right=0.3cm of sel, font=\small] {$dname='Sales'$};
\end{tikzpicture}
\end{center}

\textbf{Βελτιστοποιημένο Query Tree (Push Selection Down):}

\begin{center}
\begin{tikzpicture}[
    node distance=1.2cm,
    op/.style={circle, draw=secondary, fill=box-blue, minimum size=0.8cm},
    rel/.style={rectangle, draw=primary, fill=box-green, rounded corners, minimum width=1.5cm}]

    \node[op] (proj) {$\pi$};
    \node[op, below=of proj] (join) {$\bowtie$};
    \node[rel, below left=1cm and 0.5cm of join] (emp) {Employee};
    \node[op, below right=1cm and 0.5cm of join] (sel) {$\sigma$};
    \node[rel, below=of sel] (dept) {Department};

    \draw[arrow] (join) -- (proj);
    \draw[arrow] (emp) -- (join);
    \draw[arrow] (sel) -- (join);
    \draw[arrow] (dept) -- (sel);

    \node[right=0.3cm of proj, font=\small] {$name$};
    \node[right=0.3cm of sel, font=\small] {$dname='Sales'$};
\end{tikzpicture}
\end{center}

\begin{infobox}[Κανόνες Βελτιστοποίησης]
\begin{enumerate}
    \item \textbf{Push selections down}: Μειώνει το μέγεθος ενδιάμεσων αποτελεσμάτων
    \item \textbf{Push projections down}: Μειώνει το πλάτος tuples
    \item \textbf{Join ordering}: Μικρότερα ενδιάμεσα αποτελέσματα πρώτα
\end{enumerate}
\end{infobox}

% -----------------------------------------------------------------------------
% Exercise 4.4
% -----------------------------------------------------------------------------
\subsection{Άσκηση 4.4: Cost Comparison}

\begin{formulabox}[Εκφώνηση]
Για τον πίνακα $R$ με 10,000 tuples και index στο attribute $A$:

Συγκρίνετε το cost για: $\sigma_{A=5}(R)$

\begin{enumerate}
    \item Table scan
    \item B+-tree index (height = 3, clustering factor = 1.0)
    \item Hash index
\end{enumerate}
\end{formulabox}

\subsubsection*{Λύση}

Έστω $B_R = 100$ pages, selectivity = 1\% (100 tuples match).

\textbf{(a) Table Scan:}
$$Cost = B_R = 100 \text{ page I/Os}$$

\textbf{(b) B+-tree Index (clustered):}
$$Cost = height + \text{matching pages} = 3 + 1 = 4 \text{ I/Os}$$
(Clustered = τα matching tuples είναι συνεχόμενα)

\textbf{(c) Hash Index:}
$$Cost = 1.2 \text{ (average bucket access)} = \approx 2 \text{ I/Os}$$

\begin{center}
\begin{tabular}{lr}
\toprule
\textbf{Method} & \textbf{I/O Cost} \\
\midrule
Table Scan & 100 \\
B+-tree (clustered) & 4 \\
Hash Index & $\sim$2 \\
\bottomrule
\end{tabular}
\end{center}

\begin{warningbox}[Unclustered Index]
Αν ο B+-tree index ήταν unclustered, το cost θα ήταν:
$$Cost = 3 + 100 = 103 \text{ I/Os (χειρότερο από scan!)}$$
\end{warningbox}


% =============================================================================
% SECTION 5: Data Preprocessing
% =============================================================================
\section{Προεπεξεργασία Δεδομένων}

% -----------------------------------------------------------------------------
% Exercise 5.1
% -----------------------------------------------------------------------------
\subsection{Άσκηση 5.1: Ταξινόμηση Attribute Types (NOIR)}

\begin{formulabox}[Εκφώνηση]
Ταξινομήστε τα παρακάτω attributes σε Nominal, Ordinal, Interval, ή Ratio:

\begin{enumerate}
    \item Αριθμός τηλεφώνου
    \item Θερμοκρασία σε Celsius
    \item Εκπαιδευτικό επίπεδο (Λύκειο, Πτυχίο, Μεταπτυχιακό, PhD)
    \item Ηλικία σε έτη
    \item Χρώμα ματιών
    \item Ημερομηνία γέννησης
    \item Βαθμολογία 1-5 αστέρια
\end{enumerate}
\end{formulabox}

\subsubsection*{Λύση}

\begin{center}
\begin{tabular}{llp{7cm}}
\toprule
\textbf{Attribute} & \textbf{Type} & \textbf{Εξήγηση} \\
\midrule
Αρ. τηλεφώνου & \textbf{Nominal} & Αναγνωριστικό, δεν έχει νόημα η σειρά/αριθμητικές πράξεις \\
Θερμοκρασία (°C) & \textbf{Interval} & Διαφορές έχουν νόημα, αλλά 0°C δεν = «καμία θερμοκρασία» \\
Εκπαιδευτικό επίπεδο & \textbf{Ordinal} & Σειρά έχει νόημα, αλλά διαφορές όχι \\
Ηλικία & \textbf{Ratio} & Υπάρχει απόλυτο 0, λόγοι έχουν νόημα (διπλάσια ηλικία) \\
Χρώμα ματιών & \textbf{Nominal} & Κατηγορικό, χωρίς σειρά \\
Ημερομηνία γέννησης & \textbf{Interval} & Διαφορές έχουν νόημα, δεν υπάρχει «0 ημερομηνία» \\
Rating 1-5 & \textbf{Ordinal} & Σειρά έχει νόημα, αλλά 4★ δεν = 2×2★ \\
\bottomrule
\end{tabular}
\end{center}

\begin{infobox}[NOIR Summary]
\begin{itemize}
    \item \textbf{N}ominal: Κατηγορίες, mode μόνο
    \item \textbf{O}rdinal: + Σειρά, median
    \item \textbf{I}nterval: + Διαφορές, mean
    \item \textbf{R}atio: + Λόγοι, geometric mean
\end{itemize}
\end{infobox}

% -----------------------------------------------------------------------------
% Exercise 5.2
% -----------------------------------------------------------------------------
\subsection{Άσκηση 5.2: Equal-Width Binning}

\begin{formulabox}[Εκφώνηση]
Δίνονται οι τιμές: $\{4, 8, 9, 15, 21, 21, 24, 25, 26, 28, 29, 34\}$

Εφαρμόστε equal-width binning με 3 bins.
\end{formulabox}

\subsubsection*{Λύση}

\textbf{Βήμα 1: Υπολογισμός εύρους}

$$\text{Range} = \max - \min = 34 - 4 = 30$$

\textbf{Βήμα 2: Υπολογισμός πλάτους bin}

$$\text{Width} = \frac{\text{Range}}{\text{num\_bins}} = \frac{30}{3} = 10$$

\textbf{Βήμα 3: Καθορισμός ορίων}

\begin{itemize}
    \item Bin 1: $[4, 14)$
    \item Bin 2: $[14, 24)$
    \item Bin 3: $[24, 34]$
\end{itemize}

\textbf{Βήμα 4: Κατανομή τιμών}

\begin{center}
\begin{tabular}{cl}
\toprule
\textbf{Bin} & \textbf{Τιμές} \\
\midrule
Bin 1 $[4, 14)$ & 4, 8, 9 \\
Bin 2 $[14, 24)$ & 15, 21, 21 \\
Bin 3 $[24, 34]$ & 24, 25, 26, 28, 29, 34 \\
\bottomrule
\end{tabular}
\end{center}

\begin{warningbox}[Μειονέκτημα Equal-Width]
Τα bins μπορεί να έχουν πολύ άνιση κατανομή (Bin 3 έχει 6 τιμές, Bin 1 έχει 3).
\end{warningbox}

% -----------------------------------------------------------------------------
% Exercise 5.3
% -----------------------------------------------------------------------------
\subsection{Άσκηση 5.3: Equal-Frequency Binning}

\begin{formulabox}[Εκφώνηση]
Για τις ίδιες τιμές: $\{4, 8, 9, 15, 21, 21, 24, 25, 26, 28, 29, 34\}$

Εφαρμόστε equal-frequency (equal-depth) binning με 3 bins.
\end{formulabox}

\subsubsection*{Λύση}

\textbf{Βήμα 1: Υπολογισμός μεγέθους bin}

$$\text{Bin size} = \frac{n}{\text{num\_bins}} = \frac{12}{3} = 4 \text{ τιμές/bin}$$

\textbf{Βήμα 2: Κατανομή (sorted)}

\begin{center}
\begin{tabular}{cl}
\toprule
\textbf{Bin} & \textbf{Τιμές} \\
\midrule
Bin 1 & 4, 8, 9, 15 \\
Bin 2 & 21, 21, 24, 25 \\
Bin 3 & 26, 28, 29, 34 \\
\bottomrule
\end{tabular}
\end{center}

\begin{infobox}[Equal-Frequency vs Equal-Width]
\begin{itemize}
    \item \textbf{Equal-Width}: Ίδιο εύρος ανά bin, διαφορετικό πλήθος
    \item \textbf{Equal-Frequency}: Ίδιο πλήθος ανά bin, διαφορετικό εύρος
\end{itemize}
\end{infobox}

% -----------------------------------------------------------------------------
% Exercise 5.4
% -----------------------------------------------------------------------------
\subsection{Άσκηση 5.4: Smoothing by Bin Means}

\begin{formulabox}[Εκφώνηση]
Χρησιμοποιώντας τα equal-frequency bins από την προηγούμενη άσκηση, εφαρμόστε
smoothing by bin means.
\end{formulabox}

\subsubsection*{Λύση}

\textbf{Βήμα 1: Υπολογισμός μέσου όρου κάθε bin}

\begin{align*}
\text{Mean}_1 &= \frac{4 + 8 + 9 + 15}{4} = \frac{36}{4} = 9 \\
\text{Mean}_2 &= \frac{21 + 21 + 24 + 25}{4} = \frac{91}{4} = 22.75 \\
\text{Mean}_3 &= \frac{26 + 28 + 29 + 34}{4} = \frac{117}{4} = 29.25
\end{align*}

\textbf{Βήμα 2: Αντικατάσταση}

\begin{center}
\begin{tabular}{ccc}
\toprule
\textbf{Bin} & \textbf{Original} & \textbf{Smoothed} \\
\midrule
Bin 1 & 4, 8, 9, 15 & 9, 9, 9, 9 \\
Bin 2 & 21, 21, 24, 25 & 22.75, 22.75, 22.75, 22.75 \\
Bin 3 & 26, 28, 29, 34 & 29.25, 29.25, 29.25, 29.25 \\
\bottomrule
\end{tabular}
\end{center}

% -----------------------------------------------------------------------------
% Exercise 5.5
% -----------------------------------------------------------------------------
\subsection{Άσκηση 5.5: Smoothing by Bin Boundaries}

\begin{formulabox}[Εκφώνηση]
Για τα ίδια bins, εφαρμόστε smoothing by bin boundaries.
\end{formulabox}

\subsubsection*{Λύση}

\textbf{Μέθοδος:} Κάθε τιμή αντικαθίσταται με το πλησιέστερο όριο του bin της.

\textbf{Bin 1: [4, 15]}
\begin{itemize}
    \item 4 $\to$ 4 (είναι το κάτω όριο)
    \item 8 $\to$ 4 (πιο κοντά στο 4 παρά στο 15)
    \item 9 $\to$ 4 (πιο κοντά στο 4)
    \item 15 $\to$ 15 (είναι το άνω όριο)
\end{itemize}

\textbf{Bin 2: [21, 25]}
\begin{itemize}
    \item 21 $\to$ 21
    \item 21 $\to$ 21
    \item 24 $\to$ 25 (πιο κοντά στο 25)
    \item 25 $\to$ 25
\end{itemize}

\textbf{Bin 3: [26, 34]}
\begin{itemize}
    \item 26 $\to$ 26
    \item 28 $\to$ 26 (πιο κοντά στο 26)
    \item 29 $\to$ 26 (ή 34, ίση απόσταση - επιλέγουμε το κοντινότερο)
    \item 34 $\to$ 34
\end{itemize}

\begin{center}
\begin{tabular}{ccc}
\toprule
\textbf{Bin} & \textbf{Original} & \textbf{Smoothed} \\
\midrule
Bin 1 & 4, 8, 9, 15 & 4, 4, 4, 15 \\
Bin 2 & 21, 21, 24, 25 & 21, 21, 25, 25 \\
Bin 3 & 26, 28, 29, 34 & 26, 26, 26, 34 \\
\bottomrule
\end{tabular}
\end{center}

% -----------------------------------------------------------------------------
% Exercise 5.6
% -----------------------------------------------------------------------------
\subsection{Άσκηση 5.6: Min-Max Normalization}

\begin{formulabox}[Εκφώνηση]
Κανονικοποιήστε τις τιμές $\{10, 20, 30, 40, 50\}$ στο διάστημα $[0, 1]$
χρησιμοποιώντας Min-Max normalization.
\end{formulabox}

\subsubsection*{Λύση}

\textbf{Τύπος Min-Max Normalization:}

$$x' = \frac{x - \min}{\max - \min}$$

\textbf{Υπολογισμοί:}

$\min = 10$, $\max = 50$, $\text{range} = 40$

\begin{align*}
10' &= \frac{10 - 10}{40} = 0.00 \\
20' &= \frac{20 - 10}{40} = 0.25 \\
30' &= \frac{30 - 10}{40} = 0.50 \\
40' &= \frac{40 - 10}{40} = 0.75 \\
50' &= \frac{50 - 10}{40} = 1.00
\end{align*}

\begin{infobox}[Γενική Μορφή]
Για normalization στο διάστημα $[a, b]$:
$$x' = \frac{x - \min}{\max - \min} \cdot (b - a) + a$$
\end{infobox}

% -----------------------------------------------------------------------------
% Exercise 5.7
% -----------------------------------------------------------------------------
\subsection{Άσκηση 5.7: Z-Score Standardization}

\begin{formulabox}[Εκφώνηση]
Για τις ίδιες τιμές $\{10, 20, 30, 40, 50\}$, εφαρμόστε Z-score standardization.
\end{formulabox}

\subsubsection*{Λύση}

\textbf{Τύπος Z-Score:}

$$z = \frac{x - \mu}{\sigma}$$

\textbf{Βήμα 1: Υπολογισμός μέσου όρου}

$$\mu = \frac{10 + 20 + 30 + 40 + 50}{5} = 30$$

\textbf{Βήμα 2: Υπολογισμός τυπικής απόκλισης}

$$\sigma = \sqrt{\frac{\sum(x_i - \mu)^2}{n}} = \sqrt{\frac{400 + 100 + 0 + 100 + 400}{5}} = \sqrt{200} \approx 14.14$$

\textbf{Βήμα 3: Υπολογισμός z-scores}

\begin{align*}
z_{10} &= \frac{10 - 30}{14.14} \approx -1.41 \\
z_{20} &= \frac{20 - 30}{14.14} \approx -0.71 \\
z_{30} &= \frac{30 - 30}{14.14} = 0.00 \\
z_{40} &= \frac{40 - 30}{14.14} \approx +0.71 \\
z_{50} &= \frac{50 - 30}{14.14} \approx +1.41
\end{align*}

\begin{infobox}[Min-Max vs Z-Score]
\begin{itemize}
    \item \textbf{Min-Max}: Bounded $[0,1]$, ευαίσθητο σε outliers
    \item \textbf{Z-Score}: Unbounded, μέσος=0, std=1, λιγότερο ευαίσθητο σε outliers
\end{itemize}
\end{infobox}

% -----------------------------------------------------------------------------
% Exercise 5.8
% -----------------------------------------------------------------------------
\subsection{Άσκηση 5.8: Υπολογισμός Αποστάσεων}

\begin{formulabox}[Εκφώνηση]
Δίνονται τα διανύσματα:
\begin{itemize}
    \item $\mathbf{x} = (1, 2, 3)$
    \item $\mathbf{y} = (4, 6, 3)$
\end{itemize}

Υπολογίστε: (a) Euclidean distance, (b) Manhattan distance, (c) Cosine similarity
\end{formulabox}

\subsubsection*{Λύση}

\textbf{(a) Euclidean Distance:}

$$d_E(\mathbf{x}, \mathbf{y}) = \sqrt{\sum_{i=1}^{n} (x_i - y_i)^2}$$

$$d_E = \sqrt{(1-4)^2 + (2-6)^2 + (3-3)^2} = \sqrt{9 + 16 + 0} = \sqrt{25} = 5$$

\textbf{(b) Manhattan Distance:}

$$d_M(\mathbf{x}, \mathbf{y}) = \sum_{i=1}^{n} |x_i - y_i|$$

$$d_M = |1-4| + |2-6| + |3-3| = 3 + 4 + 0 = 7$$

\textbf{(c) Cosine Similarity:}

$$\cos(\mathbf{x}, \mathbf{y}) = \frac{\mathbf{x} \cdot \mathbf{y}}{||\mathbf{x}|| \cdot ||\mathbf{y}||}$$

\begin{align*}
\mathbf{x} \cdot \mathbf{y} &= 1 \cdot 4 + 2 \cdot 6 + 3 \cdot 3 = 4 + 12 + 9 = 25 \\
||\mathbf{x}|| &= \sqrt{1^2 + 2^2 + 3^2} = \sqrt{14} \\
||\mathbf{y}|| &= \sqrt{4^2 + 6^2 + 3^2} = \sqrt{61} \\
\cos(\mathbf{x}, \mathbf{y}) &= \frac{25}{\sqrt{14} \cdot \sqrt{61}} = \frac{25}{\sqrt{854}} \approx 0.855
\end{align*}

\begin{center}
\begin{tabular}{lr}
\toprule
\textbf{Distance Metric} & \textbf{Value} \\
\midrule
Euclidean & 5.00 \\
Manhattan & 7 \\
Cosine Similarity & 0.855 \\
Cosine Distance & $1 - 0.855 = 0.145$ \\
\bottomrule
\end{tabular}
\end{center}


% =============================================================================
% SECTION 6: Data Streams
% =============================================================================
\section{Data Streams}

% -----------------------------------------------------------------------------
% Exercise 6.1
% -----------------------------------------------------------------------------
\subsection{Άσκηση 6.1: Bloom Filter Insert \& Lookup}

\begin{formulabox}[Εκφώνηση]
Ένα Bloom Filter έχει μέγεθος $m = 10$ bits και χρησιμοποιεί 2 hash functions:
\begin{itemize}
    \item $h_1(x) = x \mod 10$
    \item $h_2(x) = (2x + 3) \mod 10$
\end{itemize}

\begin{enumerate}
    \item Εισάγετε τα στοιχεία: 15, 23, 36
    \item Ελέγξτε αν υπάρχουν: 23, 42
\end{enumerate}
\end{formulabox}

\subsubsection*{Λύση}

\textbf{Βήμα 1: Αρχικοποίηση}

Bit array: $[0, 0, 0, 0, 0, 0, 0, 0, 0, 0]$

\textbf{Βήμα 2: Insert 15}

\begin{align*}
h_1(15) &= 15 \mod 10 = 5 \\
h_2(15) &= (30 + 3) \mod 10 = 3
\end{align*}

Set bits 3, 5: $[0, 0, 0, \mathbf{1}, 0, \mathbf{1}, 0, 0, 0, 0]$

\textbf{Βήμα 3: Insert 23}

\begin{align*}
h_1(23) &= 23 \mod 10 = 3 \\
h_2(23) &= (46 + 3) \mod 10 = 9
\end{align*}

Set bits 3, 9: $[0, 0, 0, \mathbf{1}, 0, 1, 0, 0, 0, \mathbf{1}]$

\textbf{Βήμα 4: Insert 36}

\begin{align*}
h_1(36) &= 36 \mod 10 = 6 \\
h_2(36) &= (72 + 3) \mod 10 = 5
\end{align*}

Set bits 5, 6: $[0, 0, 0, 1, 0, \mathbf{1}, \mathbf{1}, 0, 0, 1]$

\textbf{Final bit array:} $[0, 0, 0, 1, 0, 1, 1, 0, 0, 1]$

\textbf{Lookup 23:}
\begin{itemize}
    \item $h_1(23) = 3$ $\to$ bit[3] = 1 $\checkmark$
    \item $h_2(23) = 9$ $\to$ bit[9] = 1 $\checkmark$
\end{itemize}
Αποτέλεσμα: \textbf{Probably in set}

\textbf{Lookup 42:}
\begin{align*}
h_1(42) &= 42 \mod 10 = 2 \\
h_2(42) &= (84 + 3) \mod 10 = 7
\end{align*}
\begin{itemize}
    \item bit[2] = 0 $\times$
\end{itemize}
Αποτέλεσμα: \textbf{Definitely not in set}

% -----------------------------------------------------------------------------
% Exercise 6.2
% -----------------------------------------------------------------------------
\subsection{Άσκηση 6.2: Bloom Filter False Positive}

\begin{formulabox}[Εκφώνηση]
Με το ίδιο Bloom Filter από την προηγούμενη άσκηση, ελέγξτε αν το 13 είναι
false positive.
\end{formulabox}

\subsubsection*{Λύση}

Bit array: $[0, 0, 0, 1, 0, 1, 1, 0, 0, 1]$

\textbf{Lookup 13:}
\begin{align*}
h_1(13) &= 13 \mod 10 = 3 \\
h_2(13) &= (26 + 3) \mod 10 = 9
\end{align*}

\begin{itemize}
    \item bit[3] = 1 $\checkmark$
    \item bit[9] = 1 $\checkmark$
\end{itemize}

Αποτέλεσμα: «Probably in set»

\textbf{Αλλά:} Το 13 \textbf{δεν} εισήχθη ποτέ! Αυτό είναι \textbf{FALSE POSITIVE}.

\begin{warningbox}[False Positive Rate]
Η πιθανότητα false positive είναι περίπου:
$$P_{fp} \approx \left(1 - e^{-kn/m}\right)^k$$
όπου $k$ = hash functions, $n$ = inserted elements, $m$ = bits.

Για $k=2$, $n=3$, $m=10$:
$$P_{fp} \approx (1 - e^{-0.6})^2 \approx (0.45)^2 \approx 0.20 = 20\%$$
\end{warningbox}

% -----------------------------------------------------------------------------
% Exercise 6.3
% -----------------------------------------------------------------------------
\subsection{Άσκηση 6.3: Count-Min Sketch Query}

\begin{formulabox}[Εκφώνηση]
Ένα Count-Min Sketch έχει $d = 2$ hash functions και $w = 5$ counters ανά row.
Μετά από updates, ο πίνακας είναι:

\begin{center}
\begin{tabular}{c|ccccc}
 & 0 & 1 & 2 & 3 & 4 \\
\hline
$h_1$ & 3 & 7 & 2 & 5 & 1 \\
$h_2$ & 4 & 2 & 8 & 3 & 6 \\
\end{tabular}
\end{center}

Αν $h_1(\text{``apple''}) = 1$ και $h_2(\text{``apple''}) = 3$, ποια είναι η εκτίμηση
για το count του ``apple''?
\end{formulabox}

\subsubsection*{Λύση}

\textbf{Count-Min Sketch Query:}

$$\hat{f}(x) = \min_{i=1}^{d} \text{CMS}[i][h_i(x)]$$

\textbf{Υπολογισμός:}
\begin{align*}
\text{CMS}[1][h_1(\text{``apple''})] &= \text{CMS}[1][1] = 7 \\
\text{CMS}[2][h_2(\text{``apple''})] &= \text{CMS}[2][3] = 3
\end{align*}

$$\hat{f}(\text{``apple''}) = \min(7, 3) = 3$$

\begin{infobox}[Γιατί Minimum?]
Τα counts μπορούν να αυξηθούν λόγω collisions με άλλα στοιχεία.
Το minimum δίνει την καλύτερη (μικρότερη) εκτίμηση, που είναι εγγυημένα $\geq$ actual count.
\end{infobox}

% -----------------------------------------------------------------------------
% Exercise 6.4
% -----------------------------------------------------------------------------
\subsection{Άσκηση 6.4: Count-Min Sketch Update}

\begin{formulabox}[Εκφώνηση]
Χρησιμοποιώντας τον ίδιο CMS, εκτελέστε: Update(``banana'', +2)

Δίνεται: $h_1(\text{``banana''}) = 4$, $h_2(\text{``banana''}) = 2$
\end{formulabox}

\subsubsection*{Λύση}

\textbf{Update operation:}

Για κάθε row $i$: $\text{CMS}[i][h_i(x)] \mathrel{+}= c$

\textbf{Πριν:}
\begin{center}
\begin{tabular}{c|ccccc}
 & 0 & 1 & 2 & 3 & 4 \\
\hline
$h_1$ & 3 & 7 & 2 & 5 & 1 \\
$h_2$ & 4 & 2 & 8 & 3 & 6 \\
\end{tabular}
\end{center}

\textbf{Update:}
\begin{align*}
\text{CMS}[1][4] &= 1 + 2 = 3 \\
\text{CMS}[2][2] &= 8 + 2 = 10
\end{align*}

\textbf{Μετά:}
\begin{center}
\begin{tabular}{c|ccccc}
 & 0 & 1 & 2 & 3 & 4 \\
\hline
$h_1$ & 3 & 7 & 2 & 5 & \textbf{3} \\
$h_2$ & 4 & 2 & \textbf{10} & 3 & 6 \\
\end{tabular}
\end{center}

% -----------------------------------------------------------------------------
% Exercise 6.5
% -----------------------------------------------------------------------------
\subsection{Άσκηση 6.5: Error Bounds}

\begin{formulabox}[Εκφώνηση]
Σχεδιάστε Count-Min Sketch με error $\varepsilon = 0.01$ και failure probability
$\delta = 0.05$. Υπολογίστε τις διαστάσεις.
\end{formulabox}

\subsubsection*{Λύση}

\textbf{Τύποι:}

$$w = \lceil e / \varepsilon \rceil \approx \lceil 2.718 / 0.01 \rceil = 272$$

$$d = \lceil \ln(1/\delta) \rceil = \lceil \ln(20) \rceil = \lceil 3.0 \rceil = 3$$

\textbf{Εγγύηση:}

Με πιθανότητα τουλάχιστον $1 - \delta = 95\%$:
$$\hat{f}(x) \leq f(x) + \varepsilon \cdot N$$

όπου $N$ = συνολικό count όλων των updates.

\textbf{Απαιτούμενη μνήμη:}
$$\text{Memory} = w \times d \times \text{counter\_size} = 272 \times 3 \times 4 = 3264 \text{ bytes}$$

% -----------------------------------------------------------------------------
% Exercise 6.6
% -----------------------------------------------------------------------------
\subsection{Άσκηση 6.6: Bloom Filter vs Count-Min Sketch}

\begin{formulabox}[Εκφώνηση]
Για κάθε σενάριο, επιλέξτε την κατάλληλη δομή (Bloom Filter ή Count-Min Sketch):

\begin{enumerate}
    \item Έλεγχος αν ένα URL έχει επισκεφθεί πριν
    \item Καταμέτρηση επισκέψεων ανά URL
    \item Spam filter: έλεγχος αν email είναι σε blacklist
    \item Εύρεση heavy hitters σε network traffic
\end{enumerate}
\end{formulabox}

\subsubsection*{Λύση}

\begin{center}
\begin{tabular}{clp{6cm}}
\toprule
\textbf{\#} & \textbf{Δομή} & \textbf{Αιτιολόγηση} \\
\midrule
1 & \textbf{Bloom Filter} & Membership test μόνο (yes/no) \\
2 & \textbf{Count-Min Sketch} & Χρειάζεται καταμέτρηση, όχι μόνο membership \\
3 & \textbf{Bloom Filter} & Membership test σε blacklist \\
4 & \textbf{Count-Min Sketch} & Heavy hitters = υψηλά counts \\
\bottomrule
\end{tabular}
\end{center}

\begin{infobox}[Σύγκριση]
\begin{itemize}
    \item \textbf{Bloom Filter}: Membership queries ($\in$?), no deletions, no counts
    \item \textbf{Count-Min Sketch}: Frequency queries (count?), supports negative updates
\end{itemize}
\end{infobox}


% =============================================================================
% SECTION 7: HMM & Sequences
% =============================================================================
\section{HMM και Διακριτές Ακολουθίες}

% -----------------------------------------------------------------------------
% Exercise 7.1
% -----------------------------------------------------------------------------
\subsection{Άσκηση 7.1: Υποακολουθίες}

\begin{formulabox}[Εκφώνηση]
Δίνεται η ακολουθία $S = \langle A, B, C, D \rangle$.

\begin{enumerate}
    \item Πόσες υποακολουθίες (subsequences) έχει;
    \item Απαριθμήστε όλες τις υποακολουθίες μήκους 2.
    \item Είναι η $\langle B, D \rangle$ υποακολουθία της $S$;
    \item Είναι η $\langle C, B \rangle$ υποακολουθία της $S$;
\end{enumerate}
\end{formulabox}

\subsubsection*{Λύση}

\textbf{(1) Αριθμός υποακολουθιών:}

Κάθε στοιχείο μπορεί να είναι ή να μην είναι στην υποακολουθία:
$$\text{Total} = 2^n = 2^4 = 16 \text{ (συμπεριλαμβανομένης της κενής)}$$

\textbf{(2) Υποακολουθίες μήκους 2:}

$$\binom{4}{2} = 6 \text{ υποακολουθίες}$$

$\langle A, B \rangle$, $\langle A, C \rangle$, $\langle A, D \rangle$,
$\langle B, C \rangle$, $\langle B, D \rangle$, $\langle C, D \rangle$

\textbf{(3) $\langle B, D \rangle$ υποακολουθία;} \textbf{ΝΑΙ}

Το B εμφανίζεται πριν το D στην $S$: $A, \mathbf{B}, C, \mathbf{D}$

\textbf{(4) $\langle C, B \rangle$ υποακολουθία;} \textbf{ΟΧΙ}

Στην $S$ το B έρχεται πριν το C, όχι μετά.

\begin{infobox}[Υποακολουθία vs Υποσυμβολοσειρά]
\begin{itemize}
    \item \textbf{Subsequence}: Διατήρηση σειράς, όχι απαραίτητα συνεχόμενα
    \item \textbf{Substring}: Συνεχόμενα στοιχεία
\end{itemize}
\end{infobox}

% -----------------------------------------------------------------------------
% Exercise 7.2
% -----------------------------------------------------------------------------
\subsection{Άσκηση 7.2: Forward Algorithm - Απλό}

\begin{formulabox}[Εκφώνηση]
Ένα HMM έχει 2 καταστάσεις (H: Hot, C: Cold) και 2 παρατηρήσεις (S: Sun, R: Rain).

\textbf{Παράμετροι:}
\begin{itemize}
    \item Initial: $\pi_H = 0.6$, $\pi_C = 0.4$
    \item Transitions: $a_{HH} = 0.7$, $a_{HC} = 0.3$, $a_{CH} = 0.4$, $a_{CC} = 0.6$
    \item Emissions: $b_H(S) = 0.8$, $b_H(R) = 0.2$, $b_C(S) = 0.3$, $b_C(R) = 0.7$
\end{itemize}

Υπολογίστε την πιθανότητα της παρατήρησης $O = \langle S, R \rangle$.
\end{formulabox}

\subsubsection*{Λύση}

\textbf{Forward Algorithm:}

$$\alpha_t(j) = P(o_1, o_2, ..., o_t, q_t = j | \lambda)$$

\textbf{Βήμα 1: Αρχικοποίηση ($t=1$, $o_1 = S$)}

\begin{align*}
\alpha_1(H) &= \pi_H \cdot b_H(S) = 0.6 \times 0.8 = 0.48 \\
\alpha_1(C) &= \pi_C \cdot b_C(S) = 0.4 \times 0.3 = 0.12
\end{align*}

\textbf{Βήμα 2: Induction ($t=2$, $o_2 = R$)}

\begin{align*}
\alpha_2(H) &= \left[\alpha_1(H) \cdot a_{HH} + \alpha_1(C) \cdot a_{CH}\right] \cdot b_H(R) \\
&= [0.48 \times 0.7 + 0.12 \times 0.4] \times 0.2 \\
&= [0.336 + 0.048] \times 0.2 = 0.384 \times 0.2 = 0.0768 \\[5pt]
\alpha_2(C) &= \left[\alpha_1(H) \cdot a_{HC} + \alpha_1(C) \cdot a_{CC}\right] \cdot b_C(R) \\
&= [0.48 \times 0.3 + 0.12 \times 0.6] \times 0.7 \\
&= [0.144 + 0.072] \times 0.7 = 0.216 \times 0.7 = 0.1512
\end{align*}

\textbf{Βήμα 3: Termination}

$$P(O | \lambda) = \alpha_2(H) + \alpha_2(C) = 0.0768 + 0.1512 = 0.228$$

\textbf{Τελική Απάντηση:} $P(\langle S, R \rangle) = 0.228$

% -----------------------------------------------------------------------------
% Exercise 7.3
% -----------------------------------------------------------------------------
\subsection{Άσκηση 7.3: Forward Algorithm - Πίνακας}

\begin{formulabox}[Εκφώνηση]
Με το ίδιο HMM, υπολογίστε $P(O = \langle S, S, R \rangle)$ χρησιμοποιώντας πίνακα.
\end{formulabox}

\subsubsection*{Λύση}

\begin{center}
\begin{tabular}{c|ccc}
\toprule
 & $t=1$ (S) & $t=2$ (S) & $t=3$ (R) \\
\midrule
$\alpha_t(H)$ & $0.6 \times 0.8 = 0.48$ & $0.2976$ & $0.0655$ \\
$\alpha_t(C)$ & $0.4 \times 0.3 = 0.12$ & $0.0648$ & $0.1197$ \\
\bottomrule
\end{tabular}
\end{center}

\textbf{Υπολογισμοί $t=2$:}
\begin{align*}
\alpha_2(H) &= [0.48 \times 0.7 + 0.12 \times 0.4] \times 0.8 = 0.384 \times 0.8 = 0.3072 \\
\alpha_2(C) &= [0.48 \times 0.3 + 0.12 \times 0.6] \times 0.3 = 0.216 \times 0.3 = 0.0648
\end{align*}

\textbf{Υπολογισμοί $t=3$:}
\begin{align*}
\alpha_3(H) &= [0.3072 \times 0.7 + 0.0648 \times 0.4] \times 0.2 = 0.2409 \times 0.2 = 0.0482 \\
\alpha_3(C) &= [0.3072 \times 0.3 + 0.0648 \times 0.6] \times 0.7 = 0.1310 \times 0.7 = 0.0917
\end{align*}

$$P(O) = 0.0482 + 0.0917 = 0.1399$$

% -----------------------------------------------------------------------------
% Exercise 7.4
% -----------------------------------------------------------------------------
\subsection{Άσκηση 7.4: Viterbi Algorithm}

\begin{formulabox}[Εκφώνηση]
Με το ίδιο HMM, βρείτε την πιο πιθανή ακολουθία καταστάσεων για $O = \langle S, R \rangle$.
\end{formulabox}

\subsubsection*{Λύση}

\textbf{Viterbi Algorithm:}

$$\delta_t(j) = \max_{q_1,...,q_{t-1}} P(q_1,...,q_{t-1}, q_t=j, o_1,...,o_t | \lambda)$$

\textbf{Βήμα 1: Αρχικοποίηση ($t=1$, $o_1 = S$)}

\begin{align*}
\delta_1(H) &= \pi_H \cdot b_H(S) = 0.6 \times 0.8 = 0.48 \\
\delta_1(C) &= \pi_C \cdot b_C(S) = 0.4 \times 0.3 = 0.12
\end{align*}

$\psi_1(H) = \psi_1(C) = 0$ (αρχή)

\textbf{Βήμα 2: Recursion ($t=2$, $o_2 = R$)}

\begin{align*}
\delta_2(H) &= \max[\delta_1(H) \cdot a_{HH}, \delta_1(C) \cdot a_{CH}] \cdot b_H(R) \\
&= \max[0.48 \times 0.7, 0.12 \times 0.4] \times 0.2 \\
&= \max[0.336, 0.048] \times 0.2 = 0.336 \times 0.2 = 0.0672 \\
\psi_2(H) &= H \text{ (επειδή } 0.336 > 0.048\text{)}
\end{align*}

\begin{align*}
\delta_2(C) &= \max[\delta_1(H) \cdot a_{HC}, \delta_1(C) \cdot a_{CC}] \cdot b_C(R) \\
&= \max[0.48 \times 0.3, 0.12 \times 0.6] \times 0.7 \\
&= \max[0.144, 0.072] \times 0.7 = 0.144 \times 0.7 = 0.1008 \\
\psi_2(C) &= H \text{ (επειδή } 0.144 > 0.072\text{)}
\end{align*}

\textbf{Βήμα 3: Termination}

$$q_2^* = \arg\max[\delta_2(H), \delta_2(C)] = \arg\max[0.0672, 0.1008] = C$$

\textbf{Βήμα 4: Backtracking}

$$q_1^* = \psi_2(C) = H$$

\textbf{Τελική Απάντηση:} Πιο πιθανή ακολουθία: $\langle H, C \rangle$

\begin{center}
\begin{tikzpicture}[
    state/.style={circle, draw=secondary, fill=box-blue, minimum size=1cm, font=\bfseries},
    node distance=3cm]

    \node[state] (h1) {H};
    \node[state, below=1.5cm of h1] (c1) {C};
    \node[state, right=of h1] (h2) {H};
    \node[state, below=1.5cm of h2] (c2) {C};

    \draw[arrow, thick, accent] (h1) -- (c2);
    \draw[arrow, dashed, gray] (h1) -- (h2);
    \draw[arrow, dashed, gray] (c1) -- (h2);
    \draw[arrow, dashed, gray] (c1) -- (c2);

    \node[above=0.3cm of h1] {$t=1$: S};
    \node[above=0.3cm of h2] {$t=2$: R};
\end{tikzpicture}
\end{center}

% -----------------------------------------------------------------------------
% Exercise 7.5
% -----------------------------------------------------------------------------
\subsection{Άσκηση 7.5: Ερμηνεία HMM Παραμέτρων}

\begin{formulabox}[Εκφώνηση]
Δίνεται HMM για αναγνώριση δραστηριότητας με καταστάσεις \{Walk, Run, Sit\} και
παρατηρήσεις \{Low, Medium, High\} (accelerometer readings).

Οι παράμετροι είναι:
\begin{center}
\begin{tabular}{c|ccc}
$A$ & Walk & Run & Sit \\
\hline
Walk & 0.7 & 0.2 & 0.1 \\
Run & 0.1 & 0.8 & 0.1 \\
Sit & 0.2 & 0.1 & 0.7 \\
\end{tabular}
\quad
\begin{tabular}{c|ccc}
$B$ & Low & Med & High \\
\hline
Walk & 0.1 & 0.7 & 0.2 \\
Run & 0.05 & 0.15 & 0.8 \\
Sit & 0.8 & 0.15 & 0.05 \\
\end{tabular}
\end{center}

Τι μας λένε αυτές οι παράμετροι για τη συμπεριφορά του μοντέλου;
\end{formulabox}

\subsubsection*{Λύση}

\textbf{Transition Matrix $A$ - Ερμηνεία:}

\begin{itemize}
    \item Υψηλές διαγώνιες τιμές (0.7-0.8): Οι δραστηριότητες τείνουν να \textbf{συνεχίζονται}
    \item $a_{Walk \to Run} = 0.2$: Σχετικά πιθανή η μετάβαση από περπάτημα σε τρέξιμο
    \item $a_{Run \to Sit} = 0.1$: Σπάνια απότομη διακοπή τρεξίματος
    \item $a_{Sit \to Run} = 0.1$: Σπάνια άμεση μετάβαση από κάθισμα σε τρέξιμο
\end{itemize}

\textbf{Emission Matrix $B$ - Ερμηνεία:}

\begin{itemize}
    \item \textbf{Sit}: Κυρίως χαμηλές ενδείξεις (0.8) - λογικό, καθιστή θέση
    \item \textbf{Run}: Κυρίως υψηλές ενδείξεις (0.8) - έντονη κίνηση
    \item \textbf{Walk}: Κυρίως μεσαίες ενδείξεις (0.7) - μέτρια κίνηση
\end{itemize}

\begin{infobox}[Χρησιμότητα HMM]
Το HMM μπορεί να «εξομαλύνει» θορυβώδεις μετρήσεις χρησιμοποιώντας την πληροφορία
ότι οι καταστάσεις τείνουν να παραμένουν σταθερές.
\end{infobox}


% =============================================================================
% SECTION 8: OLAP & Data Warehousing
% =============================================================================
\section{OLAP και Data Warehousing}

% -----------------------------------------------------------------------------
% Exercise 8.1
% -----------------------------------------------------------------------------
\subsection{Άσκηση 8.1: OLAP Operations Ταξινόμηση}

\begin{formulabox}[Εκφώνηση]
Ταξινομήστε κάθε λειτουργία σε Roll-up, Drill-down, Slice, Dice, ή Pivot:

\begin{enumerate}
    \item Από πωλήσεις ανά ημέρα σε πωλήσεις ανά μήνα
    \item Εμφάνιση μόνο των πωλήσεων του 2024
    \item Από πωλήσεις ανά χώρα σε πωλήσεις ανά πόλη
    \item Πωλήσεις για συγκεκριμένη περιοχή ΚΑΙ κατηγορία προϊόντος
    \item Αλλαγή από Product×Time σε Region×Time
\end{enumerate}
\end{formulabox}

\subsubsection*{Λύση}

\begin{center}
\begin{tabular}{clp{6cm}}
\toprule
\textbf{\#} & \textbf{Operation} & \textbf{Εξήγηση} \\
\midrule
1 & \textbf{Roll-up} & Aggregation σε υψηλότερο επίπεδο ιεραρχίας (ημέρα → μήνας) \\
2 & \textbf{Slice} & Επιλογή μιας τιμής σε μία διάσταση (year = 2024) \\
3 & \textbf{Drill-down} & Μετάβαση σε χαμηλότερο επίπεδο λεπτομέρειας \\
4 & \textbf{Dice} & Επιλογή εύρους τιμών σε πολλές διαστάσεις \\
5 & \textbf{Pivot} & Αναδιάταξη διαστάσεων/αλλαγή οπτικής γωνίας \\
\bottomrule
\end{tabular}
\end{center}

\begin{center}
\begin{tikzpicture}[scale=0.7]
    % Hierarchy visualization
    \node[draw, fill=box-blue, minimum width=3cm] at (0,3) {Year};
    \node[draw, fill=box-blue, minimum width=3cm] at (0,2) {Quarter};
    \node[draw, fill=box-blue, minimum width=3cm] at (0,1) {Month};
    \node[draw, fill=box-blue, minimum width=3cm] at (0,0) {Day};

    \draw[arrow, thick, success] (1.8,0.5) -- (1.8,2.5) node[midway, right] {Roll-up};
    \draw[arrow, thick, accent] (-1.8,2.5) -- (-1.8,0.5) node[midway, left] {Drill-down};
\end{tikzpicture}
\end{center}

% -----------------------------------------------------------------------------
% Exercise 8.2
% -----------------------------------------------------------------------------
\subsection{Άσκηση 8.2: Roll-up και Drill-down Sequence}

\begin{formulabox}[Εκφώνηση]
Δίνεται cube με διαστάσεις: Time (Day → Month → Quarter → Year),
Product (Item → Category → Department), Location (City → Region → Country).

Ξεκινώντας από (Day, Item, City), περιγράψτε τη σειρά operations για να φτάσετε σε
(Year, Department, Country).
\end{formulabox}

\subsubsection*{Λύση}

\textbf{Απαιτούμενα Roll-ups:}

\begin{enumerate}
    \item \textbf{Time}: Day → Month → Quarter → Year (3 roll-ups)
    \item \textbf{Product}: Item → Category → Department (2 roll-ups)
    \item \textbf{Location}: City → Region → Country (2 roll-ups)
\end{enumerate}

\textbf{Σύνολο:} 7 consecutive roll-up operations

\begin{center}
\begin{tikzpicture}[node distance=0.8cm, font=\small]
    \node[draw, fill=box-green, rounded corners] (start) {(Day, Item, City)};
    \node[draw, fill=box-blue, rounded corners, below=of start] (s1) {(Month, Item, City)};
    \node[draw, fill=box-blue, rounded corners, below=of s1] (s2) {(Quarter, Category, Region)};
    \node[draw, fill=box-green, rounded corners, below=of s2] (end) {(Year, Department, Country)};

    \draw[arrow] (start) -- (s1) node[midway, right] {roll-up Time};
    \draw[arrow] (s1) -- (s2) node[midway, right] {roll-up all};
    \draw[arrow] (s2) -- (end) node[midway, right] {roll-up all};
\end{tikzpicture}
\end{center}

\begin{infobox}[Optimization]
Στην πράξη, τα roll-ups μπορούν να γίνουν παράλληλα σε διαφορετικές διαστάσεις
αν υπάρχουν pre-computed aggregates.
\end{infobox}

% -----------------------------------------------------------------------------
% Exercise 8.3
% -----------------------------------------------------------------------------
\subsection{Άσκηση 8.3: Slice vs Dice}

\begin{formulabox}[Εκφώνηση]
Για cube πωλήσεων με διαστάσεις (Product, Location, Time), εκφράστε σε OLAP
operations:

\begin{enumerate}
    \item Πωλήσεις για τον Ιανουάριο 2024
    \item Πωλήσεις για Electronics σε Athens τον Ιανουάριο 2024
    \item Πωλήσεις για όλα τα προϊόντα, στην Ελλάδα, για Q1 2024
\end{enumerate}
\end{formulabox}

\subsubsection*{Λύση}

\textbf{(1) Slice}

$$\text{SLICE}(Cube, Time = \text{``Jan 2024''})$$

Μειώνει τον cube κατά μία διάσταση.

\textbf{(2) Dice}

$$\text{DICE}(Cube, Product = \text{``Electronics''}, Location = \text{``Athens''}, Time = \text{``Jan 2024''})$$

Επιλογή σε πολλαπλές διαστάσεις ταυτόχρονα.

\textbf{(3) Slice + Roll-up (ή Dice)}

\begin{enumerate}
    \item Roll-up Location: City → Country (Ελλάδα)
    \item Roll-up Time: Month → Quarter (Q1)
    \item Slice: Country = ``Greece'', Quarter = ``Q1 2024''
\end{enumerate}

% -----------------------------------------------------------------------------
% Exercise 8.4
% -----------------------------------------------------------------------------
\subsection{Άσκηση 8.4: Cube Calculations}

\begin{formulabox}[Εκφώνηση]
Ένας cube έχει 3 διαστάσεις: Product (100 values), Location (50 values),
Time (365 days).

\begin{enumerate}
    \item Πόσα cells έχει ο base cube;
    \item Πόσα συνολικά cells έχει ο πλήρης CUBE (με όλα τα aggregations);
    \item Αν κάθε cell χρειάζεται 8 bytes, πόση μνήμη απαιτείται;
\end{enumerate}
\end{formulabox}

\subsubsection*{Λύση}

\textbf{(1) Base Cube Cells:}

$$|Product| \times |Location| \times |Time| = 100 \times 50 \times 365 = 1,825,000$$

\textbf{(2) Full CUBE:}

Για κάθε διάσταση, μπορεί να έχουμε τις τιμές της ή το ALL (aggregation):

$$(|Product| + 1) \times (|Location| + 1) \times (|Time| + 1)$$
$$= 101 \times 51 \times 366 = 1,884,666$$

Εναλλακτικά (για cuboids):
$$\text{Cuboids} = 2^3 = 8$$

\begin{center}
\begin{tabular}{lc}
\toprule
\textbf{Cuboid} & \textbf{Cells} \\
\midrule
(P, L, T) & $100 \times 50 \times 365 = 1,825,000$ \\
(P, L, *) & $100 \times 50 = 5,000$ \\
(P, *, T) & $100 \times 365 = 36,500$ \\
(*, L, T) & $50 \times 365 = 18,250$ \\
(P, *, *) & $100$ \\
(*, L, *) & $50$ \\
(*, *, T) & $365$ \\
(*, *, *) & $1$ \\
\midrule
\textbf{Total} & \textbf{1,885,266} \\
\bottomrule
\end{tabular}
\end{center}

\textbf{(3) Memory:}

$$1,885,266 \times 8 \text{ bytes} = 15,082,128 \text{ bytes} \approx 14.4 \text{ MB}$$

% -----------------------------------------------------------------------------
% Exercise 8.5
% -----------------------------------------------------------------------------
\subsection{Άσκηση 8.5: Star Schema Design}

\begin{formulabox}[Εκφώνηση]
Σχεδιάστε star schema για ηλεκτρονικό κατάστημα με τα εξής requirements:
\begin{itemize}
    \item Παρακολούθηση πωλήσεων (ποσότητα, τιμή, έκπτωση)
    \item Διαστάσεις: Προϊόν, Πελάτης, Ημερομηνία, Κατάστημα
    \item Ιεραρχία προϊόντων: Product → Subcategory → Category → Department
\end{itemize}
\end{formulabox}

\subsubsection*{Λύση}

\begin{center}
\begin{tikzpicture}[
    fact/.style={rectangle, draw=accent, fill=box-red, minimum width=3cm, minimum height=2cm, font=\small, align=center},
    dim/.style={rectangle, draw=secondary, fill=box-blue, minimum width=2.5cm, minimum height=1.8cm, font=\small, align=center}]

    % Fact table
    \node[fact] (fact) {
        \textbf{Sales\_Fact}\\
        \underline{product\_sk}\\
        \underline{customer\_sk}\\
        \underline{date\_sk}\\
        \underline{store\_sk}\\
        quantity\\
        unit\_price\\
        discount
    };

    % Dimension tables
    \node[dim, above left=1cm and 2cm of fact] (product) {
        \textbf{Dim\_Product}\\
        \underline{product\_sk}\\
        product\_id\\
        name\\
        subcategory\\
        category\\
        department
    };

    \node[dim, above right=1cm and 2cm of fact] (customer) {
        \textbf{Dim\_Customer}\\
        \underline{customer\_sk}\\
        customer\_id\\
        name\\
        city\\
        country
    };

    \node[dim, below left=1cm and 2cm of fact] (date) {
        \textbf{Dim\_Date}\\
        \underline{date\_sk}\\
        date\\
        day\_of\_week\\
        month\\
        quarter\\
        year
    };

    \node[dim, below right=1cm and 2cm of fact] (store) {
        \textbf{Dim\_Store}\\
        \underline{store\_sk}\\
        store\_id\\
        name\\
        city\\
        region
    };

    % Connections
    \draw[thick] (product) -- (fact);
    \draw[thick] (customer) -- (fact);
    \draw[thick] (date) -- (fact);
    \draw[thick] (store) -- (fact);
\end{tikzpicture}
\end{center}

\begin{infobox}[Surrogate Keys (\_sk)]
Χρησιμοποιούμε surrogate keys (τεχνητά κλειδιά) αντί για natural keys για:
\begin{itemize}
    \item Καλύτερη απόδοση (μικρότερα integers)
    \item Διαχείριση slowly changing dimensions
    \item Ανεξαρτησία από source systems
\end{itemize}
\end{infobox}

% -----------------------------------------------------------------------------
% Exercise 8.6
% -----------------------------------------------------------------------------
\subsection{Άσκηση 8.6: OLAP Queries σε SQL}

\begin{formulabox}[Εκφώνηση]
Χρησιμοποιώντας το star schema της προηγούμενης άσκησης, γράψτε SQL queries για:

\begin{enumerate}
    \item Συνολικές πωλήσεις ανά κατηγορία προϊόντος
    \item Top 5 πελάτες με βάση το συνολικό ποσό αγορών
    \item Πωλήσεις ανά μήνα για το 2024
\end{enumerate}
\end{formulabox}

\subsubsection*{Λύση}

\textbf{(1) Πωλήσεις ανά κατηγορία (Roll-up):}

\begin{lstlisting}[language=SQL]
SELECT
    p.category,
    SUM(f.quantity * f.unit_price * (1 - f.discount)) AS total_revenue
FROM Sales_Fact f
JOIN Dim_Product p ON f.product_sk = p.product_sk
GROUP BY p.category
ORDER BY total_revenue DESC;
\end{lstlisting}

\textbf{(2) Top 5 πελάτες:}

\begin{lstlisting}[language=SQL]
SELECT
    c.name,
    SUM(f.quantity * f.unit_price * (1 - f.discount)) AS total_spent
FROM Sales_Fact f
JOIN Dim_Customer c ON f.customer_sk = c.customer_sk
GROUP BY c.customer_sk, c.name
ORDER BY total_spent DESC
LIMIT 5;
\end{lstlisting}

\textbf{(3) Πωλήσεις ανά μήνα 2024 (Slice + Roll-up):}

\begin{lstlisting}[language=SQL]
SELECT
    d.year,
    d.month,
    SUM(f.quantity * f.unit_price * (1 - f.discount)) AS monthly_revenue
FROM Sales_Fact f
JOIN Dim_Date d ON f.date_sk = d.date_sk
WHERE d.year = 2024
GROUP BY d.year, d.month
ORDER BY d.month;
\end{lstlisting}

\begin{warningbox}[Performance]
Στα data warehouses, οι dimension tables είναι συνήθως denormalized για αποφυγή
πολλαπλών JOINs. Αυτό αυξάνει το μέγεθος αλλά βελτιώνει την απόδοση των queries.
\end{warningbox}

